\centerline{\bf \Huge Алгебраический разнобой}

\begin{flushright}
        \scriptsize{Если я сказал не брал, значит не отдам\\ Урахаре Киске}
\end{flushright}

\problem{1}
Два различных числа $x$ и $y$ (не обязательно целых) таковы, что

$$
x^2-2000 x=y^2-2000 y .
$$


Найдите сумму чисел $x$ и $y$.

\problem{2}
Саша придумал два целых числа и перемножил их. Затем взял числа, обратные придуманным, и сложил. Полученные сумма и произведение оказались обратными числами. Найдите разность модулей Сашиных чисел.

\problem{3}
50 бизнесменов - японцы, корейцы и китайцы - сидят за круглым столом. Известно, что между каждыми двумя ближайшими японцами сидит ровно столько китайцев, сколько всего за столом корейцев. Сколько китайцев может быть за столом?


\problem{4}
Вычислите разность сумм $(1 \cdot 3+3 \cdot 5+5 \cdot 7+\ldots+99 \cdot 101)-\left(1^2+3^2+5^2+\ldots+99^2\right)$.

\problem{5}
Чему равно $\left(1+\frac{1}{2}\right)\left(1+\frac{1}{3}\right)\left(1+\frac{1}{4}\right) \ldots\left(1+\frac{1}{100}\right)$ ?


\problem{6}
Найдите значение выражения $\left(1-\frac{1}{4}\right)\left(1-\frac{1}{9}\right)\left(1-\frac{1}{16}\right) \ldots\left(1-\frac{1}{100^2}\right)$.


\problem{7}
Найдите сумму $\frac{2^2}{1 \cdot 3}+\frac{4^2}{3 \cdot 5}+\frac{6^2}{5 \cdot 7}+\ldots+\frac{100^2}{99 \cdot 101}$.